\documentclass{beamer}
\usepackage[utf8]{inputenc}
\usepackage[T1]{fontenc}
\usepackage{graphicx}
\usepackage{amsmath}
\usepackage{amssymb}
\usepackage{listings}
\usepackage{xcolor}
\usepackage{tikz}
\usepackage{booktabs}

% Theme configuration
\usetheme{Madrid}
\usecolortheme{default}

% Color configuration
\definecolor{codegreen}{rgb}{0,0.6,0}
\definecolor{codegray}{rgb}{0.5,0.5,0.5}
\definecolor{codepurple}{rgb}{0.58,0,0.82}
\definecolor{backcolour}{rgb}{0.95,0.95,0.92}

% Python code style
\lstdefinestyle{pythonstyle}{
    backgroundcolor=\color{backcolour},   
    commentstyle=\color{codegreen},
    keywordstyle=\color{blue},
    numberstyle=\tiny\color{codegray},
    stringstyle=\color{codepurple},
    basicstyle=\ttfamily\footnotesize,
    breakatwhitespace=false,         
    breaklines=true,                 
    captionpos=b,                    
    keepspaces=true,                 
    numbers=left,                    
    numbersep=5pt,                  
    showspaces=false,                
    showstringspaces=false,
    showtabs=false,                  
    tabsize=2,
    language=Python
}

\lstset{style=pythonstyle}

% Document information
\title{Regression in Machine Learning}
\subtitle{Complete Course}
\author{Meyssem Bouzaien}
\institute{Academic Year: 2025 - 2026}
\date{}

\begin{document}

% Title page
\frame{\titlepage}

% Outline
\begin{frame}
\frametitle{Course Outline}
\tableofcontents
\end{frame}

% ==========================================
% SECTION 1: INTRODUCTION
% ==========================================
\section{What is Regression?}

\begin{frame}
\frametitle{Classification vs Regression}
\begin{center}
\includegraphics[width=\textwidth]{images/01_classification_vs_regression.png}
\end{center}
\end{frame}

\begin{frame}
\frametitle{Classification vs Regression - Differences}
\begin{columns}
\begin{column}{0.48\textwidth}
\textbf{Classification}
\begin{itemize}
    \item \textbf{Output:} Category (class)
    \item \textbf{Type:} Discrete variable
    \item \textbf{Examples:}
    \begin{itemize}
        \item Spam or Not spam
        \item Cat, Dog, Bird
        \item Sick or Healthy
    \end{itemize}
\end{itemize}
\end{column}
\begin{column}{0.48\textwidth}
\textbf{Regression}
\begin{itemize}
    \item \textbf{Output:} Numeric value
    \item \textbf{Type:} Continuous variable
    \item \textbf{Examples:}
    \begin{itemize}
        \item House price
        \item Temperature
        \item Salary
    \end{itemize}
\end{itemize}
\end{column}
\end{columns}

\vspace{1em}
\begin{block}{Simple Rule}
Predicting a number → Regression | Predicting a category → Classification
\end{block}
\end{frame}

\begin{frame}
\frametitle{Definition of Regression}
\begin{block}{Definition}
Regression is a supervised Machine Learning technique that predicts a continuous numeric value from input variables (features).
\end{block}

\vspace{1em}
\textbf{Principle:}
\begin{itemize}
    \item Find the relationship between the features (X) and the target (Y)
    \item Model this relationship mathematically
    \item Use this model to make predictions
\end{itemize}

\vspace{1em}
\textbf{Objective:}
$$\text{Minimize the error between predictions and actual values}$$
\end{frame}

\begin{frame}
\frametitle{Concrete Examples of Regression}
\vspace{-0.2cm}
\begin{center}
\includegraphics[width= 0.8\textwidth]{images/02_exemples_regression.png}
\end{center}
\end{frame}

\begin{frame}
\frametitle{Applications of Regression}
\begin{table}
\centering
\small
\begin{tabular}{lll}
\toprule
\textbf{Domain} & \textbf{Input (X)} & \textbf{Output (Y)} \\
\midrule
Real estate & Surface, location & Price \\
Finance & Stock history & Future price \\
Marketing & Advertising budget & Sales \\
Weather & Historical data & Temperature \\
Healthcare & Age, weight, blood pressure & Disease risk \\
E-commerce & Season, promotions & Product demand \\
\bottomrule
\end{tabular}
\end{table}

\begin{alertblock}{Important}
Regression is everywhere in industry!
\end{alertblock}
\end{frame}

% ==========================================
% SECTION 2: SIMPLE LINEAR REGRESSION
% ==========================================
\section{Simple Linear Regression}

\begin{frame}
\frametitle{Principle of Linear Regression}
\begin{center}
\includegraphics[width=0.95\textwidth]{images/03_principe_regression_lineaire.png}
\end{center}
\end{frame}

\begin{frame}
\frametitle{What is Linear Regression?}
\begin{block}{Definition}
Linear regression models the relationship between an input variable X and an output variable Y using a straight line.
\end{block}

\vspace{0.5em}
\textbf{Idea:}
\begin{itemize}
    \item Find THE best line that runs "through the middle" of the points
    \item This line represents the general trend
    \item We then use it to predict new values
\end{itemize}

\vspace{0.5em}
\textbf{Why ``linear''?}
\begin{itemize}
    \item Because we're looking for a straight-line relationship
    \item Simple, fast, interpretable
\end{itemize}
\end{frame}

\begin{frame}
\frametitle{Equation of a Line}
\begin{center}
\includegraphics[width=0.8\textwidth]{images/04_equation_droite.png}
\end{center}
\end{frame}

\begin{frame}
\frametitle{Mathematical Equation}
\begin{block}{Equation of the line}
$$y = ax + b$$
\end{block}

\vspace{1em}
\begin{itemize}
    \item \textbf{$y$}: Variable to predict (price, temperature, salary...)
    \item \textbf{$x$}: Input variable (surface, experience, mileage...)
    \item \textbf{$a$}: Slope (coefficient) - indicates the effect of X on Y
    \item \textbf{$b$}: Intercept - value of Y when X=0
\end{itemize}

\vspace{1em}
\begin{exampleblock}{Example: Real Estate Price}
$$\text{Price} = 2000 \times \text{Surface} + 50000$$
\begin{itemize}
    \item $a = 2000$: each m² adds €2000
    \item $b = 50000$: base price
\end{itemize}
\end{exampleblock}
\end{frame}

\begin{frame}
\frametitle{Residuals and Errors}
\begin{center}
\includegraphics[width=\textwidth]{images/05_residus_mse.png}
\end{center}
\end{frame}

\begin{frame}
\frametitle{How Do We Measure Quality?}
\begin{block}{Residual}
Residual = Actual value - Predicted value
$$e_i = y_i - \hat{y}_i$$
\end{block}

\vspace{0.5em}
\textbf{Mean Squared Error (MSE):}
$$\text{MSE} = \frac{1}{n}\sum_{i=1}^{n}(y_i - \hat{y}_i)^2$$

\vspace{0.5em}
\textbf{Goal of regression:}
\begin{itemize}
    \item Find the parameters $a$ and $b$ that minimize the MSE
    \item The smaller the MSE, the better the model
\end{itemize}

\vspace{0.5em}
\begin{alertblock}{Attention}
MSE strongly penalizes large errors (because they're squared)!
\end{alertblock}
\end{frame}

\begin{frame}[fragile]
\frametitle{Simple Linear Regression - Python Code}
\begin{lstlisting}[language=Python]
from sklearn.linear_model import LinearRegression
import numpy as np

# Data
X = np.array([50, 60, 70, 80, 90, 100]).reshape(-1, 1)
y = np.array([150, 180, 210, 240, 270, 300])  # Price (k$)

# Create the model
model = LinearRegression()
# Train
model.fit(X, y)
# Coefficients
print(f"Slope (a): {model.coef_[0]}")
print(f"Intercept (b): {model.intercept_}")
# Predict
new_surface = np.array([[75]])
predicted_price = model.predict(new_surface)
print(f"Price for 75m2: {predicted_price[0]}k$")
\end{lstlisting}
\end{frame}

\begin{frame}
\frametitle{Complete Example: Real Estate Price}
\vspace{-0.2cm}
\begin{center}
\includegraphics[width=0.8\textwidth]{images/06_exemple_complet_immobilier.png}
\end{center}
\end{frame}

% ==========================================
% SECTION 3: MULTIPLE LINEAR REGRESSION
% ==========================================
\section{Multiple Linear Regression}

\begin{frame}
\frametitle{Simple vs Multiple Regression}
\begin{center}
\includegraphics[width=\textwidth]{images/07_regression_multiple.png}
\end{center}
\end{frame}

\begin{frame}
\frametitle{What is Multiple Regression?}
\begin{block}{Definition}
Multiple linear regression uses several input variables to predict Y.
\end{block}

\vspace{0.5em}
\textbf{Equation:}
$$y = a_1 x_1 + a_2 x_2 + a_3 x_3 + ... + a_n x_n + b$$

\vspace{0.5em}
\begin{exampleblock}{Example: Real Estate Price}
$$\text{Price} = 1500 \times \text{Surface} + 20000 \times \text{Bedrooms} - 500 \times \text{Age} + 50000$$
\begin{itemize}
    \item Several combined features
    \item Each feature has its own coefficient
    \item More accurate but more complex model
\end{itemize}
\end{exampleblock}
\end{frame}

\begin{frame}
\frametitle{Advantages of Multiple Regression}
\textbf{Why use several features?}

\vspace{1em}
\begin{itemize}
    \item \textbf{More accurate:} Captures more information
    \item \textbf{More realistic:} The real world is multidimensional
    \item \textbf{Control of variables:} Isolates the effect of each feature
\end{itemize}

\vspace{1em}
\begin{exampleblock}{Example}
To predict a car's price:
\begin{itemize}
    \item Mileage
    \item Year
    \item Brand
    \item Horsepower
    \item Condition
\end{itemize}
→ More features = more accurate prediction
\end{exampleblock}
\end{frame}

\begin{frame}[fragile]
\frametitle{Multiple Regression - Python Code}
\begin{lstlisting}[language=Python]
from sklearn.linear_model import LinearRegression
import pandas as pd
# Data with several features
data = {
    'Surface': [50, 60, 70, 80, 90, 100],
    'Bedrooms': [1, 2, 2, 3, 3, 4],
    'Age': [5, 10, 3, 15, 8, 2],
    'Price': [150, 180, 220, 230, 280, 320]
}
df = pd.DataFrame(data)
# Features and target
X = df[['Surface', 'Bedrooms', 'Age']]
y = df['Price']
# Model
model = LinearRegression()
model.fit(X, y)
# Coefficients
for feature, coef in zip(X.columns, model.coef_):
    print(f"{feature}: {coef:.2f}")
print(f"Intercept: {model.intercept_:.2f}")
\end{lstlisting}
\end{frame}

% ==========================================
% SECTION 4: POLYNOMIAL REGRESSION
% ==========================================
\section{Polynomial Regression}

\begin{frame}
\frametitle{Beyond the Straight Line}
\begin{block}{Question}
What if the relationship between X and Y isn't linear?
\end{block}

\vspace{1em}
\textbf{Polynomial Regression:}
\begin{itemize}
    \item Uses powers of X: $x$, $x^2$, $x^3$, etc.
    \item Allows modeling curves
    \item More flexible than simple linear regression
\end{itemize}

\vspace{1em}
\textbf{General equation (degree 2):}
$$y = a_2 x^2 + a_1 x + b$$

\vspace{0.5em}
\textbf{Example:}
$$\text{Sales} = -2 \times \text{Temp}^2 + 80 \times \text{Temp} - 500$$
\end{frame}

\begin{frame}
\frametitle{Different Polynomial Degrees}
\vspace{-0.2cm}

\begin{center}
\includegraphics[width=0.8\textwidth]{images/08_regression_polynomiale_degres.png}
\end{center}
\end{frame}

\begin{frame}
\frametitle{Choosing the Right Degree}
\begin{table}
\centering
\begin{tabular}{lll}
\toprule
\textbf{Degree} & \textbf{Shape} & \textbf{Use} \\
\midrule
1 & Straight line & Linear relationship \\
2 & Parabola & Simple curve (U or ∩) \\
3-5 & Complex curves & Non-linear relationships \\
>10 & Very complex & Risk of overfitting! \\
\bottomrule
\end{tabular}
\end{table}

\vspace{1em}
\begin{alertblock}{Watch the Degree!}
\begin{itemize}
    \item Degree too small → Underfitting (too simple)
    \item Degree too large → Overfitting (too complex)
    \item Always test on held-out data!
\end{itemize}
\end{alertblock}
\end{frame}

\begin{frame}[fragile]
\frametitle{Polynomial Regression - Python Code}
\begin{lstlisting}[language=Python]
from sklearn.preprocessing import PolynomialFeatures
from sklearn.linear_model import LinearRegression
import numpy as np
# Data
X = np.array([1, 2, 3, 4, 5]).reshape(-1, 1)
y = np.array([2, 8, 18, 32, 50])  # y = 2*x^2
# Create polynomial features (degree 2)
poly = PolynomialFeatures(degree=2)
X_poly = poly.fit_transform(X)
# X_poly now contains: [1, x, x^2]
# Regression
model = LinearRegression()
model.fit(X_poly, y)
# Predict
X_new = np.array([[6]])
X_new_poly = poly.transform(X_new)
y_pred = model.predict(X_new_poly)
print(f"Prediction for x=6: {y_pred[0]}")
\end{lstlisting}
\end{frame}

% ==========================================
% SECTION 5: EVALUATION METRICS
% ==========================================
\section{Evaluation Metrics}

\begin{frame}
\frametitle{The 4 Main Metrics}
\vspace{-0.2cm}

\begin{center}
\includegraphics[width=0.8\textwidth]{images/09_metriques_evaluation.png}
\end{center}
\end{frame}

\begin{frame}
\frametitle{MSE - Mean Squared Error}
\begin{block}{Formula}
$$\text{MSE} = \frac{1}{n}\sum_{i=1}^{n}(y_i - \hat{y}_i)^2$$
\end{block}

\textbf{Characteristics:}
\begin{itemize}
    \item Average of squared errors
    \item Strongly penalizes large errors
    \item Always positive
    \item Unit = (unit of Y)²
\end{itemize}

\vspace{0.5em}
\begin{exampleblock}{Example}
If Y is in euros and MSE = 1,000,000, this means:
\begin{itemize}
    \item The squared errors average 1M€²
    \item Hard to interpret → use RMSE instead
\end{itemize}
\end{exampleblock}
\end{frame}

\begin{frame}
\frametitle{RMSE - Root Mean Squared Error}
\begin{block}{Formula}
$$\text{RMSE} = \sqrt{\text{MSE}} = \sqrt{\frac{1}{n}\sum_{i=1}^{n}(y_i - \hat{y}_i)^2}$$
\end{block}

\textbf{Characteristics:}
\begin{itemize}
    \item Square root of the MSE
    \item \textbf{Same unit as Y!} (easy to interpret)
    \item Represents the "typical" error
\end{itemize}

\vspace{0.5em}
\begin{exampleblock}{Example}
Real estate price: RMSE = €25,000
\begin{itemize}
    \item On average, predictions are off by ±€25k
    \item Directly comparable to prices
\end{itemize}
\end{exampleblock}

\textbf{When to use it?} When you want a metric in the unit of Y
\end{frame}

\begin{frame}
\frametitle{MAE - Mean Absolute Error}
\begin{block}{Formula}
$$\text{MAE} = \frac{1}{n}\sum_{i=1}^{n}|y_i - \hat{y}_i|$$
\end{block}

\textbf{Characteristics:}
\begin{itemize}
    \item Average of absolute errors
    \item Less sensitive to outliers than MSE/RMSE
    \item Same unit as Y
    \item More robust
\end{itemize}

\vspace{0.5em}
\textbf{MSE vs MAE:}
\begin{itemize}
    \item MSE/RMSE: strongly penalizes large errors
    \item MAE: treats all errors equally
\end{itemize}

\vspace{0.5em}
\textbf{When to use it?} When you don't want to over-penalize outliers
\end{frame}

\begin{frame}
\frametitle{R² - Coefficient of Determination}
\begin{block}{Formula}
$$R^2 = 1 - \frac{\sum(y_i - \hat{y}_i)^2}{\sum(y_i - \bar{y})^2}$$
\end{block}

\textbf{Characteristics:}
\begin{itemize}
    \item Ranges between 0 and 1 (sometimes negative if very poor)
    \item Measures the \% of variance explained
    \item Unitless (easy to compare)
\end{itemize}

\vspace{0.5em}
\textbf{Interpretation:}
\begin{itemize}
    \item $R^2 = 1.0$: Perfect prediction (100\%)
    \item $R^2 = 0.8$: 80\% of variance explained
    \item $R^2 = 0$: The model explains nothing
\end{itemize}

\vspace{0.5em}
\textbf{When to use it?} Standard metric for comparing models
\end{frame}

\begin{frame}
\frametitle{Interpreting R²}
\begin{center}
\includegraphics[width=\textwidth]{images/10_interpretation_r2.png}
\end{center}
\end{frame}

\begin{frame}
\frametitle{Metrics Summary Table}
\begin{table}
\centering
\small
\begin{tabular}{llll}
\toprule
\textbf{Metric} & \textbf{Unit} & \textbf{Best} & \textbf{Usage} \\
\midrule
MSE & Y² & 0 & Mathematics \\
RMSE & Y & 0 & Interpretation \\
MAE & Y & 0 & Robustness \\
R² & None & 1 & Comparison \\
\bottomrule
\end{tabular}
\end{table}

\vspace{1em}
\textbf{Practical advice:}
\begin{itemize}
    \item Use \textbf{R²} for an overall view
    \item Use \textbf{RMSE} for business interpretation
    \item Use \textbf{MAE} if there are many outliers
\end{itemize}
\end{frame}

\begin{frame}[fragile]
\frametitle{Computing the Metrics - Python Code}
\begin{lstlisting}[language=Python]
from sklearn.metrics import mean_squared_error, r2_score, mean_absolute_error
import numpy as np

# Actual and predicted values
y_true = np.array([100, 120, 90, 110, 95])
y_pred = np.array([105, 115, 88, 112, 98])

# Compute the metrics
mse = mean_squared_error(y_true, y_pred)
rmse = np.sqrt(mse)
mae = mean_absolute_error(y_true, y_pred)
r2 = r2_score(y_true, y_pred)

# Display
print(f"MSE  : {mse:.2f}")
print(f"RMSE : {rmse:.2f}")
print(f"MAE  : {mae:.2f}")
print(f"R2   : {r2:.3f}")
\end{lstlisting}
\end{frame}

% ==========================================
% SECTION 6: TRAIN/TEST AND VALIDATION
% ==========================================
\section{Train/Test Split and Validation}

\begin{frame}
\frametitle{Train/Test Split}
\begin{center}
\includegraphics[width=\textwidth]{images/11_train_test_split.png}
\end{center}
\end{frame}

\begin{frame}
\frametitle{Why Split the Data?}
\begin{block}{Golden Rule}
NEVER evaluate on the training data!
\end{block}

\vspace{1em}
\textbf{Reason:}
\begin{itemize}
    \item The model "knows" the training data
    \item It can memorize it (overfitting)
    \item Training error doesn't reflect true performance
\end{itemize}

\vspace{1em}
\textbf{Solution: Train/Test Split}
\begin{itemize}
    \item \textbf{Train Set (70-80\%)}: Train the model
    \item \textbf{Test Set (20-30\%)}: Evaluate performance
    \item The test set simulates "new data"
\end{itemize}
\end{frame}

\begin{frame}
\frametitle{Overfitting vs Underfitting}
\begin{center}
\includegraphics[width=\textwidth]{images/12_overfitting_underfitting.png}
\end{center}
\end{frame}

\begin{frame}
\frametitle{Understanding Overfitting and Underfitting}
\begin{columns}
\begin{column}{0.48\textwidth}
\textbf{Underfitting}
\begin{itemize}
    \item Model too simple
    \item Doesn't capture the patterns
    \item High bias
    \item Poor on train AND test
\end{itemize}

\vspace{0.5em}
\textbf{Solutions:}
\begin{itemize}
    \item More complex model
    \item More features
    \item Higher polynomial degree
\end{itemize}
\end{column}

\begin{column}{0.48\textwidth}
\textbf{Overfitting}
\begin{itemize}
    \item Model too complex
    \item Memorizes the data
    \item High variance
    \item Good on train, poor on test
\end{itemize}

\vspace{0.5em}
\textbf{Solutions:}
\begin{itemize}
    \item More data
    \item Regularization
    \item Simpler model
\end{itemize}
\end{column}
\end{columns}

\vspace{1em}
\begin{alertblock}{Detection}
Always compare Train vs Test performance!
\end{alertblock}
\end{frame}

\begin{frame}[fragile]
\frametitle{Train/Test Split - Python Code}
\begin{lstlisting}[language=Python]
from sklearn.model_selection import train_test_split
from sklearn.linear_model import LinearRegression
from sklearn.metrics import r2_score
# Data
X = [[1], [2], [3], [4], [5], [6], [7], [8], [9], [10]]
y = [2, 4, 6, 8, 10, 12, 14, 16, 18, 20]
# 70/30 Split
X_train, X_test, y_train, y_test = train_test_split(
    X, y, test_size=0.3, random_state=42
)
# Train
model = LinearRegression()
model.fit(X_train, y_train)
# Evaluate
y_pred_train = model.predict(X_train)
y_pred_test = model.predict(X_test)

print(f"R2 Train : {r2_score(y_train, y_pred_train):.3f}")
print(f"R2 Test  : {r2_score(y_test, y_pred_test):.3f}")
\end{lstlisting}
\end{frame}

% ==========================================
% SECTION 7: REGULARIZATION
% ==========================================
\section{Regularization: Ridge and Lasso}

\begin{frame}
\frametitle{What is Regularization?}
\begin{block}{Definition}
Regularization penalizes coefficients that are too large in order to prevent overfitting.
\end{block}

\vspace{1em}
\textbf{Problem:}
\begin{itemize}
    \item Complex models → very large coefficients
    \item Large coefficients → overfitting
\end{itemize}

\vspace{1em}
\textbf{Solution:}
\begin{itemize}
    \item Add a penalty to the cost function
    \item Force the model to use smaller coefficients
    \item Better generalization
\end{itemize}

\vspace{1em}
\textbf{Two main types:}
\begin{itemize}
    \item \textbf{Ridge (L2)}: Penalizes the square of the coefficients
    \item \textbf{Lasso (L1)}: Penalizes the absolute value
\end{itemize}
\end{frame}

\begin{frame}
\frametitle{Ridge vs Lasso}
\vspace{-0.2cm}
\begin{center}
\includegraphics[width=0.8\textwidth]{images/13_regularisation_ridge_lasso.png}
\end{center}
\end{frame}

\begin{frame}
\frametitle{Ridge Regression (L2)}
\begin{block}{Cost function}
$$\text{Cost} = \text{MSE} + \alpha \sum_{i=1}^{n} a_i^2$$
\end{block}

\vspace{0.5em}
\textbf{Characteristics:}
\begin{itemize}
    \item Penalizes the square of the coefficients
    \item Shrinks coefficients but doesn't set them to zero
    \item Keeps all features
    \item Good when many features matter
\end{itemize}

\vspace{0.5em}
\textbf{Parameter $\alpha$:}
\begin{itemize}
    \item $\alpha$ = 0: classic linear regression
    \item $\alpha$ small: little regularization
    \item $\alpha$ large: strong regularization
\end{itemize}
\end{frame}

\begin{frame}
\frametitle{Lasso Regression (L1)}
\begin{block}{Cost function}
$$\text{Cost} = \text{MSE} + \alpha \sum_{i=1}^{n} |a_i|$$
\end{block}

\vspace{0.5em}
\textbf{Characteristics:}
\begin{itemize}
    \item Penalizes the absolute value of the coefficients
    \item Can set some coefficients to exactly zero
    \item Performs automatic feature selection
    \item Good when many features are useless
\end{itemize}

\vspace{0.5em}
\textbf{Main advantage:}
\begin{itemize}
    \item Simplifies the model by removing features
    \item More interpretable
    \item Identifies the important features
\end{itemize}
\end{frame}

\begin{frame}
\frametitle{Which to Use When?}
\begin{table}
\centering
\begin{tabular}{lll}
\toprule
\textbf{Situation} & \textbf{Regularization} & \textbf{Reason} \\
\midrule
Many features & Lasso & Auto selection \\
Correlated features & Ridge & Keeps everything \\
p > n (more features than data) & Lasso or Ridge & Avoids overfitting \\
Interpretability & Lasso & Fewer features \\
Max performance & Test both & Depends on data \\
\bottomrule
\end{tabular}
\end{table}

\vspace{1em}
\begin{block}{ElasticNet}
Combines Ridge and Lasso: the best of both worlds!
\end{block}
\end{frame}

\begin{frame}[fragile]
\frametitle{Ridge and Lasso - Python Code}
\begin{lstlisting}[language=Python]
from sklearn.linear_model import Ridge, Lasso
from sklearn.preprocessing import PolynomialFeatures

# Data
X_train, y_train = ...  # your data

# Polynomial features
poly = PolynomialFeatures(degree=5)
X_poly = poly.fit_transform(X_train)
# Ridge
ridge = Ridge(alpha=10)
ridge.fit(X_poly, y_train)
# Lasso
lasso = Lasso(alpha=0.5)
lasso.fit(X_poly, y_train)
# Compare the coefficients
print("Ridge coefs:", ridge.coef_)
print("Lasso coefs:", lasso.coef_)
# Lasso will have some coefficients = 0
\end{lstlisting}
\end{frame}

% ==========================================
% SECTION 8: COMPARISON AND CHOICE
% ==========================================
\section{Comparing the Models}

\begin{frame}
\frametitle{Comparing the Different Models}
\vspace{-0.2cm}

\begin{center}
\includegraphics[width=0.8\textwidth]{images/14_comparaison_modeles.png}
\end{center}
\end{frame}

\begin{frame}
\frametitle{Comparison Table}
\begin{table}
\centering
\tiny
\begin{tabular}{lllll}
\toprule
\textbf{Model} & \textbf{Complexity} & \textbf{Interpretability} & \textbf{Advantages} & \textbf{Disadvantages} \\
\midrule
Simple linear & Low & Excellent & Simple, fast & Linear only \\
Multiple linear & Medium & Good & Several features & Linear \\
Polynomial & High & Medium & Curves & Overfitting \\
Ridge & Medium & Good & Anti-overfitting & Hyperparameter \\
Lasso & Medium & Excellent & Feature selection & Hyperparameter \\
\bottomrule
\end{tabular}
\end{table}

\vspace{1em}
\textbf{How to choose?}
\begin{enumerate}
    \item Start with simple linear regression
    \item Add features (multiple) if needed
    \item Try polynomial if non-linear
    \item Regularize (Ridge/Lasso) if overfitting
    \item Compare R² on the test set
\end{enumerate}
\end{frame}


% SECTION 9: COMPLETE WORKFLOW
% ==========================================
\section{Complete Workflow}

\begin{frame}
\frametitle{Complete Regression Pipeline}
\vspace{-0.2cm}

\begin{center}
\includegraphics[width=0.8\textwidth]{images/15_workflow_complet.png}
\end{center}
\end{frame}

\begin{frame}
\frametitle{The 9 Steps of the Workflow}
\begin{enumerate}
    \item \textbf{Raw data}: Collection and loading
    \item \textbf{Exploration}: Visualization, statistics
    \item \textbf{Correlations}: Identify relationships
    \item \textbf{Train/Test Split}: 70/30 or 80/20 split
    \item \textbf{Training}: Fit the model
    \item \textbf{Predictions}: Predict on the test set
    \item \textbf{Evaluation}: Compute metrics (R², RMSE...)
    \item \textbf{Residual analysis}: Check for patterns
    \item \textbf{Deployment}: Use in production
\end{enumerate}

\vspace{0.5em}
\begin{alertblock}{Important}
This is an iterative process! We often go back to earlier steps.
\end{alertblock}
\end{frame}



% ==========================================
% SECTION 10: BEST PRACTICES
% ==========================================



\begin{frame}
\frametitle{When to Use Which Model?}
\begin{table}
\centering
\small
\begin{tabular}{ll}
\toprule
\textbf{Situation} & \textbf{Recommended model} \\
\midrule
Simple linear relationship & Simple linear regression \\
Several features & Multiple regression \\
Curved relationship & Polynomial regression \\
Many features & Lasso (selection) \\
Correlated features & Ridge \\
Overfitting & Ridge or Lasso \\
Need for interpretability & Linear or Lasso \\
Max performance & Test several! \\
\bottomrule
\end{tabular}
\end{table}

\vspace{1em}
\begin{block}{Advice}
In practice, always test several models and compare their performance!
\end{block}
\end{frame}

% ==========================================
% SECTION 11: PRACTICE EXERCISES
% ==========================================
\section{Practice Exercises}

\begin{frame}
\frametitle{Exercise 1: Car Prices}
\textbf{Dataset:} Used cars

\vspace{0.5em}
\textbf{Features:}
\begin{itemize}
    \item Mileage
    \item Year
    \item Horsepower (HP)
    \item Number of doors
\end{itemize}

\vspace{0.5em}
\textbf{Target:} Price (in €)

\vspace{1em}
\textbf{To do:}
\begin{enumerate}
    \item Load and explore the data
    \item Visualize the correlations
    \item Build a multiple regression model
    \item Evaluate with R² and RMSE
    \item Predict the price of a car: 50000km, 2018, 120HP, 5 doors
\end{enumerate}
\end{frame}

\begin{frame}
\frametitle{Exercise 2: Sales Forecasting}
\textbf{Dataset:} Monthly sales for a store

\vspace{0.5em}
\textbf{Features:}
\begin{itemize}
    \item Advertising budget (€)
    \item Number of promotions
    \item Average temperature
    \item Public holiday (yes/no)
\end{itemize}

\vspace{0.5em}
\textbf{Target:} Number of sales

\vspace{1em}
\textbf{To do:}
\begin{enumerate}
    \item Train/Test split (80/20)
    \item Test linear and polynomial regression (degree 2)
    \item Compare the R² scores
    \item Identify the 2 most important features
    \item Predict sales with: €5000 ad spend, 3 promos, 25°C, not a holiday
\end{enumerate}
\end{frame}

\begin{frame}
\frametitle{Exercise 3: Regularization}
\textbf{Dataset:} Real estate prices with 20 features

\vspace{1em}
\textbf{Problem:} Many features, risk of overfitting

\vspace{1em}
\textbf{To do:}
\begin{enumerate}
    \item Build a linear model without regularization
    \item Note R² on train and test
    \item Build a Ridge model ($\alpha= 10)
    \item Build a Lasso model ($\alpha= 0.5)
    \item Compare the 3 models
    \item Identify the features selected by Lasso
    \item Which one generalizes best?
\end{enumerate}
\end{frame}



% ==========================================
% SECTION 12: CONCLUSION
% ==========================================
\section{Conclusion}

\begin{frame}
\frametitle{What You Have Learned}
\textbf{Fundamental concepts:}
\begin{itemize}
    \item Difference between classification/regression
    \item Simple and multiple linear regression
    \item Polynomial regression
    \item Regularization (Ridge, Lasso)
\end{itemize}

\vspace{0.5em}
\textbf{Metrics:}
\begin{itemize}
    \item MSE, RMSE, MAE, R²
    \item Interpretation and use
\end{itemize}

\vspace{0.5em}
\textbf{Validation:}
\begin{itemize}
    \item Train/Test split
    \item Overfitting vs Underfitting
    \item Residual analysis
\end{itemize}

\vspace{0.5em}
\textbf{Practice:}
\begin{itemize}
    \item Complete workflow with Scikit-learn
    \item Model comparison
\end{itemize}
\end{frame}




\begin{frame}
\frametitle{Recommended Resources}


\textbf{Datasets to practice:}
\begin{itemize}
    \item Kaggle: \texttt{kaggle.com/datasets}
    \item UCI ML Repository
    \item Scikit-learn: built-in datasets
\end{itemize}
\end{frame}

\begin{frame}
\frametitle{Key Points to Remember}
\begin{block}{Regression = Predicting numbers}
Classification → categories | Regression → continuous values
\end{block}


\begin{block}{Always Train/Test Split}
Never evaluate on the training data!
\end{block}


\begin{block}{R² is your friend}
The main metric for comparing models
\end{block}

\begin{block}{Start simple}
Linear → Multiple → Polynomial → Regularized
\end{block}

\begin{block}{Overfitting = the enemy}
Watch train vs test, regularize if needed
\end{block}
\end{frame}



\end{document}
